\documentclass[10.5pt]{article}

\usepackage[margin=0.55in, top=0.5in, bottom=0.5in]{geometry}
\usepackage[T1]{fontenc}
\usepackage[utf8]{inputenc}
\usepackage{lmodern}
\usepackage{array}
\usepackage{booktabs}
\usepackage{enumitem}
\usepackage{hyperref}
\usepackage{parskip}
\usepackage{ragged2e}
\usepackage{setspace}
\usepackage{tabularx}
\usepackage{titlesec}
\usepackage[table]{xcolor}
\usepackage{tikz}
\usepackage[most]{tcolorbox}
\usetikzlibrary{arrows.meta,positioning,shapes.geometric,calc}

\definecolor{PrimaryRed}{HTML}{8B1E1E}
\definecolor{SoftBlue}{HTML}{DDEEFF}
\definecolor{DarkGray}{HTML}{333333}
\definecolor{SoftGreen}{HTML}{EEF6F2}
\definecolor{SoftGold}{HTML}{FFF6E6}
\definecolor{SoftPurple}{HTML}{F7F2FA}
\definecolor{SoftRose}{HTML}{FCEEEF}

\hypersetup{colorlinks=true,linkcolor=black,urlcolor=black,citecolor=black}
\setlength{\parindent}{0pt}
\setstretch{1.06}
\newcolumntype{Y}{>{\RaggedRight\arraybackslash}X}
\newcolumntype{L}[1]{>{\RaggedRight\arraybackslash}p{#1}}
\setlist[itemize]{topsep=1pt,itemsep=0pt,parsep=0pt,partopsep=0pt,leftmargin=1.2em}
\setlist[enumerate]{topsep=1pt,itemsep=1pt,parsep=0pt,partopsep=0pt,leftmargin=1.3em}
\titleformat{\section}{\normalsize\bfseries\color{PrimaryRed}}{}{0em}{}[\vspace{-4pt}]
\titlespacing{\section}{0pt}{5pt}{2pt}

\newtcolorbox{ReflectionBox}{
  colback=SoftPurple, colframe=PrimaryRed, boxrule=0.7pt, arc=2mm,
  left=1.5mm, right=1.5mm, top=1mm, bottom=1mm
}
\newtcolorbox{TakeawayBox}{
  colback=SoftGreen, colframe=DarkGray, boxrule=0.5pt, arc=2mm,
  left=1.5mm, right=1.5mm, top=0.8mm, bottom=0.8mm
}

\begin{document}

% ── HEADER ──────────────────────────────────────────────────────────────
\begin{center}
{\small University of Rochester \textbf{·} Warner School of Education \textbf{·} ED452C}\\[2pt]
{\color{PrimaryRed}\rule{\textwidth}{1.2pt}}\\[3pt]
{\large\bfseries Journal Article Review: Social-Emotional Learning, Classroom Climate, and Inclusive Design}\\[1pt]
{\small\itshape An Evidence-Based Reflection on the RULER Approach to SEL}\\[2pt]
{\color{PrimaryRed}\rule{0.72\textwidth}{0.6pt}}\\[3pt]
\footnotesize
\begin{tabular}{@{}r@{\ }l@{\quad}r@{\ }l@{}}
  \textbf{Student:} & Piter Zacari Garcia Bautista &
  \textbf{Course:} & ED452C \\
  \textbf{Assignment:} & Journal Article Review --- SEL / PBIS &
  \textbf{Date:} & \today \\
\end{tabular}
\end{center}
\vspace{2pt}

% ── GUIDING CLAIM ───────────────────────────────────────────────────────
\begin{ReflectionBox}
\footnotesize\textbf{Guiding claim:} A well-implemented SEL program does not merely teach students to manage feelings---it redesigns the relational environment so that neurodivergent, CLD, chronically ill, trauma-experienced, and multiply identified students can access learning without first performing neurotypicality as the price of belonging. This is the gap the Puzzle Plan closes, and the standard EQUITAS holds every classroom system to.
\end{ReflectionBox}
\vspace{2pt}

% ── CITATION ────────────────────────────────────────────────────────────
\section{Article}
\footnotesize
Rivers, S.\ E., Brackett, M.\ A., Reyes, M.\ R., Elbertson, N.\ A., \& Salovey, P.\ (2013). Improving the social and emotional climate of classrooms: A clustered randomized controlled trial testing the RULER Approach. \textit{Prevention Science, 14}(1), 77--87. \url{https://doi.org/10.1007/s11121-012-0305-2}

% ── SUMMARY ─────────────────────────────────────────────────────────────
\section{Summary}
\footnotesize
\textbf{Purpose:} Tested whether the RULER Approach to SEL (Yale Center for Emotional Intelligence) could improve the observed social-emotional climate of upper-elementary classrooms in a large urban district (Rivers et al., 2013). RULER embeds emotional recognition, labeling, expression, and regulation into existing content---mirroring the Puzzle Plan's design-not-extract principle.

\textbf{Design \& sample:} Clustered randomized controlled trial; 62 schools randomly assigned to RULER (5th--6th grade ELA) or business-as-usual control; hierarchical linear modeling accounted for student-classroom-school nesting (Rivers et al., 2013).

\textbf{Measures:} Trained blind observers rated classroom climate on teacher--student warmth, student--student connection, student autonomy/leadership, and teacher responsiveness to student interests (Rivers et al., 2013).

\textbf{Findings:} RULER classrooms scored significantly higher on teacher--student warmth and connectedness, student autonomy and leadership opportunities, and teacher focus on student motivation over task compliance---effects large enough to detect through structured observation across a full school year (Rivers et al., 2013).

\begin{TakeawayBox}
\footnotesize\textbf{Study snapshot:} 62 urban schools · randomized design · blind observers · multi-level modeling · measurably warmer relationships, greater student autonomy, stronger teacher responsiveness (Rivers et al., 2013).
\end{TakeawayBox}
\vspace{2pt}

% ── REFLECTION ──────────────────────────────────────────────────────────
\section{What I Learned and Why It Matters}

\footnotesize
\textbf{SEL as access, not just wellness:} This article reframed SEL from a lesson-delivery model to a climate-design model. For students with ADHD, anxiety, chronic illness, or trauma histories, the emotional quality of the environment determines whether learning is even possible (Rivers et al., 2013). RULER's changes to teacher warmth and responsiveness function as Universal Design for Learning: they reduce the cognitive and emotional load students with disabilities carry into every lesson (Brackett et al., 2012; CASEL, n.d.). EQUITAS asks this question of every tool: does it create fair, multiple pathways---or demand that students perform normalcy as the price of access?

\textbf{Climate as bridge between SEL and disability supports:} When a teacher uses emotion language daily and co-creates classroom norms, a student can say ``I am overwhelmed'' without being marked as \textit{the student with the IEP}. RULER's infusion model---SEL inside ELA, not extracted from it---is a template for how differentiation can be invisible rather than stigmatizing (Rivers et al., 2013). This connects directly to IDEA's mandate for least restrictive environments and positive behavioral supports (Individuals with Disabilities Education Act, 34 C.F.R.\ \S\S\ 300.43, 300.320, 300.321).

\textbf{Implementation equity through an EQUITAS lens:} A warm climate alone does not undo structural barriers for CLD students, English language learners, or students experiencing housing instability. SEL must be paired with culturally sustaining relationship-building and material support (CASEL, n.d.; Garcia Bautista, 2025). RULER is only as equitable as the consistency and cultural responsiveness of its implementation.

\textbf{Connection to transition and adult readiness:} Students who practice autonomy, help-seeking, and co-regulation in school are better positioned to self-advocate in college and workplace settings. Transition planning must account for the relational environment, not only skills and services (Individuals with Disabilities Education Act, 34 C.F.R.\ \S\ 300.43).

\begin{ReflectionBox}
\footnotesize\textbf{Closing reflection:} The difference between a school that \textit{accommodates} disability and one that \textit{designs} for it is the difference between reactive patch-work and proactive architecture. SEL at the classroom climate level---as Rivers et al.\ (2013) demonstrate empirically---is one of the most scalable tools we have for that proactive design. Students who look like me, and lived like I did, should not have to wait years to be believed or to belong.
\end{ReflectionBox}
\vspace{3pt}

% ── REFERENCES ──────────────────────────────────────────────────────────
\section*{References}
\footnotesize\RaggedRight
Brackett, M.\ A., Rivers, S.\ E., Reyes, M.\ R., \& Salovey, P.\ (2012). Enhancing academic performance and social and emotional competence with the RULER feeling words curriculum. \textit{Learning and Individual Differences, 22}(2), 218--224. \url{https://doi.org/10.1016/j.lindif.2010.10.002}

CASEL. (n.d.). \textit{What is the CASEL framework?} \url{https://casel.org/fundamentals-of-sel/what-is-the-casel-framework/}

Garcia Bautista, P.\ Z. (2025). \textit{EQUITAS: Equity-Quantified Unbiased Intelligence Through Algorithmic Solutions---A framework for fairness-aware diagnostics in CLD populations}. Unpublished framework, Rochester Institute of Technology.

Individuals with Disabilities Education Act, 34 C.F.R.\ \S\S\ 300.43, 300.320, 300.321. \url{https://sites.ed.gov/idea/regs/}

Rivers, S.\ E., Brackett, M.\ A., Reyes, M.\ R., Elbertson, N.\ A., \& Salovey, P.\ (2013). Improving the social and emotional climate of classrooms: A clustered randomized controlled trial testing the RULER Approach. \textit{Prevention Science, 14}(1), 77--87. \url{https://doi.org/10.1007/s11121-012-0305-2}

\end{document}
